\chapter{实数的完备性}

\section{关于实数集完备性的基本定理}

在第一、二章中, 我们证明了关于实数集的确界原理和数列的单调有界定理, 给出了数列的致密性定理和柯西收敛准则. 这些命题以不同方式反映了实数集 \(\mathbf{R}\) 的一种特性, 通常称为实数的完备性或实数的连续性. 可以举例说明, 有理数集就不具有这种特性 (本节习题 4 ). 有关实数集完备性的基本定理, 除上面这些定理外, 还有区间套定理、 聚点定理和有限覆盖定理. 在本节中将阐述这三个基本定理, 并指出所有这六个基本定理的等价性.

\subsection{区间套定理}

%定义  1
\begin{definition}{}{def:}

\end{definition}
设闭区间列 \(\left\{  \left\lbrack  {{a}_{n},{b}_{n}}\right\rbrack  \right\}\) 具有如下性质:

(i) \(\left\lbrack  {{a}_{n},{b}_{n}}\right\rbrack   \supset  \left\lbrack  {{a}_{n + 1},{b}_{n + 1}}\right\rbrack  ,n = 1,2,\cdots\) ;

(ii) \(\mathop{\lim }\limits_{{n \rightarrow  \infty }}\left( {{b}_{n} - {a}_{n}}\right)  = 0\) ,

则称 \(\left\{  \left\lbrack  {{a}_{n},{b}_{n}}\right\rbrack  \right\}\) 为闭区间套,或简称区间套.

这里性质 (i) 表明, 构成区间套的闭区间列是前一个套着后一个, 即各闭区间的端点满足如下不等式:
\[
{a}_{1} \leq  {a}_{2} \leq  \cdots  \leq  {a}_{n} \leq  \cdots  \leq  {b}_{n} \leq  \cdots  \leq  {b}_{2} \leq  {b}_{1}. \tag{1}
\]
%定理 7.1
\begin{theorem}{}{the:}

\end{theorem}
 (区间套定理) 若 \(\left\{  \left\lbrack  {{a}_{n},{b}_{n}}\right\rbrack  \right\}\) 是一个区间套,则在实数系中存在唯一的一点 \(\xi\) ,使得 \(\xi  \in  \left\lbrack  {{a}_{n},{b}_{n}}\right\rbrack  ,n = 1,2,\cdots\) ,即
\[
{a}_{n} \leq  \xi  \leq  {b}_{n},n = 1,2,\cdots . \tag{2}
\]
%证
\begin{proof}

\end{proof}
由 (1) 式, \(\left\{  {a}_{n}\right\}\) 为递增有界数列,依单调有界定理, \(\left\{  {a}_{n}\right\}\) 有极限 \(\xi\) ,且有
\[
{a}_{n} \leq  \xi ,n = 1,2,\cdots \text{ . } \tag{3}
\]
同理,递减有界数列 \(\left\{  {b}_{n}\right\}\) 也有极限,并按区间套的条件(ii),有
\[
\mathop{\lim }\limits_{{n \rightarrow  \infty }}{b}_{n} = \mathop{\lim }\limits_{{n \rightarrow  \infty }}{a}_{n} = \xi , \tag{4}
\]
且.
\[
{b}_{n} \geq  \xi ,n = 1,2,\cdots \text{ . } \tag{5}
\]
联合 (3)、(5) 即得 (2) 式.

最后证明满足 (2) 的 \(\xi\) 是唯一的. 设数 \({\xi }^{\prime }\) 也满足
\[
{a}_{n} \leq  {\xi }^{\prime } \leq  {b}_{n},n = 1,2,\cdots ,
\]
则由 (2) 式有
\[
\left| {\xi  - {\xi }^{\prime }}\right|  \leq  {b}_{n} - {a}_{n},n = 1,2,\cdots .
\]
由区间套的条件 (ii) 得
\[
\left| {\xi  - {\xi }^{\prime }}\right|  \leq  \mathop{\lim }\limits_{{n \rightarrow  \infty }}\left( {{b}_{n} - {a}_{n}}\right)  = 0,
\]
故有 \({\xi }^{\prime } = \xi\) .

由 (4) 式容易推得如下很有用的区间套性质:

推论 若 \(\xi  \in  \left\lbrack  {{a}_{n},{b}_{n}}\right\rbrack  \left( {n = 1,2,\cdots }\right)\) 是区间套 \(\left\{  \left\lbrack  {{a}_{n},{b}_{n}}\right\rbrack  \right\}\) 所确定的点,则对任给的 \(\varepsilon  > 0\) ,存在 \(N > 0\) ,使得当 \(n > N\) 时,有
\[
\left\lbrack  {{a}_{n},{b}_{n}}\right\rbrack   \subset  U\left( {\xi ;\varepsilon }\right) .
\]
%注
\begin{remark}

\end{remark}
区间套定理中要求各个区间都是闭区间, 才能保证定理的结论成立. 对于开区间列,如 \(\left\{  \left( {0,\frac{1}{n}}\right) \right\}\) ,虽然其中各个开区间也是前一个包含后一个,且 \(\mathop{\lim }\limits_{{n \rightarrow  \infty }}\left( {\frac{1}{n} - 0}\right)  = 0\) , 但不存在属于所有开区间的公共点.

%例 1
\begin{example}

\end{example}
用区间套定理证明连续函数根的存在性定理.

%证
\begin{proof}

\end{proof}
设 \(f\) 在区间 \(\left\lbrack  {a,b}\right\rbrack\) 上连续, \(f\left( a\right) f\left( b\right)  < 0\) ,并且记 \(\left\lbrack  {{a}_{1},{b}_{1}}\right\rbrack   = \left\lbrack  {a,b}\right\rbrack\) . 令 \({c}_{1} = \frac{{a}_{1} + {b}_{1}}{2}\) , 如果 \(f\left( {c}_{1}\right)  = 0\) ,结论已经成立,故可设 \(f\left( {c}_{1}\right)  \neq  0\) . 那么 \(f\left( {a}_{1}\right) f\left( {c}_{1}\right)\) 与 \(f\left( {c}_{1}\right) f\left( {b}_{1}\right)\) 有一个小于零,不妨设 \(f\left( {a}_{1}\right) f\left( {c}_{1}\right)  < 0\) ,记 \(\left\lbrack  {{a}_{2},{b}_{2}}\right\rbrack   = \left\lbrack  {{a}_{1},{c}_{1}}\right\rbrack\) . 再令 \({c}_{2} = \frac{{a}_{2} + {b}_{2}}{2}\) ,如果 \(f\left( {c}_{2}\right)  = 0\) ,结论已经成立,故同样可设 \(f\left( {c}_{2}\right)  \neq  0\) . 那么 \(f\) 在 \(\left\lbrack  {{a}_{2},{c}_{2}}\right\rbrack\) 与 \(\left\lbrack  {{c}_{2},{b}_{2}}\right\rbrack\) 这两个区间中的某一个区间上端点值异号,并记这个区间为 \(\left\lbrack  {{a}_{3},{b}_{3}}\right\rbrack\) . 将这个过程无限重复下去,就得到一列闭区间 \(\left\{  \left\lbrack  {{a}_{n},{b}_{n}}\right\rbrack  \right\}\) ,满足:

(1) \(\left\lbrack  {{a}_{n},{b}_{n}}\right\rbrack   \supset  \left\lbrack  {{a}_{n + 1},{b}_{n + 1}}\right\rbrack  ,n = 1,2,\cdots\) ;

(2) \(\mathop{\lim }\limits_{{n \rightarrow  \infty }}\left( {{b}_{n} - {a}_{n}}\right)  = \mathop{\lim }\limits_{{n \rightarrow  \infty }}\frac{b - a}{{2}^{n - 1}} = 0\) ；

(3) \(f\left( {a}_{n}\right) f\left( {b}_{n}\right)  < 0,n = 1,2,\cdots\) .

由 (1) 和 (2) 可知 \(\left\{  \left\lbrack  {{a}_{n},{b}_{n}}\right\rbrack  \right\}\) 是一个区间套,由定理 7.1,存在 \(\xi  \in  \left\lbrack  {{a}_{n},{b}_{n}}\right\rbrack  ,n = 1\) , \(2,\cdots\) ,且有 \(\mathop{\lim }\limits_{{n \rightarrow  \infty }}{a}_{n} = \mathop{\lim }\limits_{{n \rightarrow  \infty }}{b}_{n} = \xi\) . 因为 \(f\) 在点 \(\xi\) 连续,所以由 (3) 得
\[
{f}^{2}\left( \xi \right)  = \mathop{\lim }\limits_{{n \rightarrow  \infty }}f\left( {a}_{n}\right) f\left( {b}_{n}\right)  \leq  0,
\]
则必有 \(f\left( \xi \right)  = 0\) . 显然 \(\xi  \in  \left\lbrack  {a,b}\right\rbrack\) ,它就是 \(f\) 的一个零点.

\subsection{聚点定理与有限覆盖定理}

%定义  2
\begin{definition}{}{def:}

\end{definition}
设 \(S\) 为数轴上的点集, \(\xi\) 为定点 (它可以属于 \(S\) ,也可以不属于 \(S\) ). 若 \(\xi\) 的任何邻域都含有 \(S\) 中无穷多个点,则称 \(\xi\) 为点集 \(S\) 的一个聚点.

例如,点集 \(S = \left\{  {{\left( -1\right) }^{n} + \frac{1}{n}}\right\}\) 有两个聚点 \({\xi }_{1} =  - 1\) 和 \({\xi }_{2} = 1\) ; 点集 \(S = \left\{  \frac{1}{n}\right\}\) 只有一个聚点 \(\xi  = 0\) ; 又若 \(S\) 为开区间(a, b),则(a, b)上每一点以及端点 \(a,b\) 都是 \(S\) 的聚点; 而正整数集 \({\mathbf{N}}_{ + }\) 没有聚点,任何有限数集也没有聚点.

聚点概念的另两个等价定义如下.

定义 \({2}^{\prime }\) 对于点集 \(S\) ,若点 \(\xi\) 的任何 \(\varepsilon\) 邻域都含有 \(S\) 中异于 \(\xi\) 的点,即 \({U}^{ \circ  }\left( {\xi ;\varepsilon }\right)  \cap  S \neq  \varnothing\) ,则称 \(\xi\) 为 \(S\) 的一个聚点.

定义 \({2}^{\prime \prime }\) 若存在各项互异的收敛数列 \(\left\{  {x}_{n}\right\}   \subset  S\) ,则其极限 \(\mathop{\lim }\limits_{{n \rightarrow  \infty }}{x}_{n} = \xi\) 称为 \(S\) 的一个聚点.

关于以上三个定义等价性的证明, 我们简述如下.

定义 \(2 \Rightarrow\) 定义 \({2}^{\prime }\) 是显然的,定义 \({2}^{\prime \prime } \Rightarrow\) 定义 2 也不难得到; 现证定义 \({2}^{\prime } \Rightarrow\) 定义 \({2}^{\prime \prime }\) :

设 \(\xi\) 为 \(S\) (按定义 \({2}^{\prime }\) ) 的聚点,则对任给的 \(\varepsilon  > 0\) ,存在 \(x \in  {U}^{ \circ  }\left( {\xi ;\varepsilon }\right)  \cap  S\) .

令 \({\varepsilon }_{1} = 1\) ,则存在 \({x}_{1} \in  {U}^{ \circ  }\left( {\xi ;{\varepsilon }_{1}}\right)  \cap  S\) ;

令 \({\varepsilon }_{2} = \min \left\{  {\frac{1}{2},\left| {\xi  - {x}_{1}}\right| }\right\}\) ,则存在 \({x}_{2} \in  {U}^{ \circ  }\left( {\xi ;{\varepsilon }_{2}}\right)  \cap  S\) ,且显然 \({x}_{2} \neq  {x}_{1}\) ;

............

令 \({\varepsilon }_{n} = \min \left\{  {\frac{1}{n},\left| {\xi  - {x}_{n - 1}}\right| }\right\}\) ,则存在 \({x}_{n} \in  {U}^{ \circ  }\left( {\xi ;{\varepsilon }_{n}}\right)  \cap  S\) ,且 \({x}_{n}\) 与 \({x}_{1},\cdots ,{x}_{n - 1}\) 互异.

无限地重复以上步骤,得到 \(S\) 中各项互异的数列 \(\left\{  {x}_{n}\right\}\) ,且由 \(\left| {\xi  - {x}_{n}}\right|  < {\varepsilon }_{n} \leq  \frac{1}{n}\) ,易见 \(\mathop{\lim }\limits_{{n \rightarrow  \infty }}{x}_{n} = \xi\)

%定理 7.2
\begin{theorem}{}{the:}

\end{theorem}
 (魏尔斯特拉斯 (Weierstrass) 聚点定理) 实轴上的任一有界无限点集 \(S\) 至少有一个聚点.

%证
\begin{proof}

\end{proof}
证法一 因 \(S\) 为有界点集,故存在 \(M > 0\) ,使得 \(S \subset  \left\lbrack  {-M,M}\right\rbrack\) ,记 \(\left\lbrack  {{a}_{1},{b}_{1}}\right\rbrack   = \left\lbrack  {-M,M}\right\rbrack\) .

现将 \(\left\lbrack  {{a}_{1},{b}_{1}}\right\rbrack\) 等分为两个子区间. 因 \(S\) 为无限点集,故两个子区间中至少有一个含有 \(S\) 中无穷多个点,记此子区间为 \(\left\lbrack  {{a}_{2},{b}_{2}}\right\rbrack\) ,则 \(\left\lbrack  {{a}_{1},{b}_{1}}\right\rbrack   \supset  \left\lbrack  {{a}_{2},{b}_{2}}\right\rbrack\) ,且
\[
{b}_{2} - {a}_{2} = \frac{1}{2}\left( {{b}_{1} - {a}_{1}}\right)  = M.
\]
再将 \(\left\lbrack  {{a}_{2},{b}_{2}}\right\rbrack\) 等分为两个子区间,则其中至少有一个子区间含有 \(S\) 中无穷多个点, 取出这样的一个子区间,记为 \(\left\lbrack  {{a}_{3},{b}_{3}}\right\rbrack\) ,则 \(\left\lbrack  {{a}_{2},{b}_{2}}\right\rbrack   \supset  \left\lbrack  {{a}_{3},{b}_{3}}\right\rbrack\) ,且
\[
{b}_{3} - {a}_{3} = \frac{1}{2}\left( {{b}_{2} - {a}_{2}}\right)  = \frac{M}{2}.
\]
将此等分子区间的步骤无限地进行下去,得到一个区间列 \(\left\{  \left\lbrack  {{a}_{n},{b}_{n}}\right\rbrack  \right\}\) ,它满足
\[
\left\lbrack  {{a}_{n},{b}_{n}}\right\rbrack   \supset  \left\lbrack  {{a}_{n + 1},{b}_{n + 1}}\right\rbrack  ,n = 1,2,\cdots ,
\]
\[
{b}_{n} - {a}_{n} = \frac{M}{{2}^{n - 2}} \rightarrow  0\left( {n \rightarrow  \infty }\right) ,
\]
即 \(\left\{  \left\lbrack  {{a}_{n},{b}_{n}}\right\rbrack  \right\}\) 是区间套,且其中每一个闭区间都含有 \(S\) 中无穷多个点.

由区间套定理,存在唯一的一点 \(\xi  \in  \left\lbrack  {{a}_{n},{b}_{n}}\right\rbrack  ,n = 1,2,\cdots\) . 于是由定理 7.1 的推论, 对任给的 \(\varepsilon  > 0\) ,存在 \(N > 0\) ,当 \(n > N\) 时有 \(\left\lbrack  {{a}_{n},{b}_{n}}\right\rbrack   \subset  U\left( {\xi ;\varepsilon }\right)\) . 从而 \(U\left( {\xi ;\varepsilon }\right)\) 内含有 \(S\) 中无穷多个点,按定义 \(2,\xi\) 为 \(S\) 的一个聚点.

证法二 设 \(S\) 是有界无限点集. 在 \(S\) 中取一列两两不同的点列 \(\left\{  {x}_{n}\right\}\) ,显然 \(\left\{  {x}_{n}\right\}\) 是有界点列. 由致密性定理, \(\left\{  {x}_{n}\right\}\) 存在一个收敛的子列 \(\left\{  {x}_{{n}_{k}}\right\}\) ,其极限设为 \({x}_{0}\) . 那么对于任意正数 \(\varepsilon\) ,存在 \(K\) ,当 \(k > K\) 时,有 \({x}_{0} - \varepsilon  < {x}_{{n}_{k}} < {x}_{0} + \varepsilon\) . 这就说明 \(\left( {{x}_{0} - \varepsilon ,{x}_{0} + \varepsilon }\right)\) 含有 \(S\) 中无限多个点,即 \({x}_{0}\) 是 \(S\) 的一个聚点.

十分明显,致密性定理是聚点定理的一种特殊情形. 这只需把有界数列 \(\left\{  {x}_{n}\right\}\) 看成有界点集 \(S\) ,并把 \(\left\{  {x}_{n}\right\}\) 中的无限多个“项”看成 \(S\) 中的无限多个“点”.

%定义  3
\begin{definition}{}{def:}

\end{definition}
设 \(S\) 为数轴上的点集, \(H\) 为开区间的集合 (即 \(H\) 的每一个元素都是形如 \(\left( {\alpha ,\beta }\right)\) 的开区间). 若 \(S\) 中任何一点都含在 \(H\) 中至少一个开区间内,则称 \(H\) 为 \(S\) 的一个开覆盖,或称 \(H\) 覆盖 \(S\) . 若 \(H\) 中开区间的个数是无限 (有限) 的,则称 \(H\) 为 \(S\) 的一个无限开覆盖 (有限开覆盖).

在具体问题中,一个点集的开覆盖常由该问题的某些条件所确定. 例如,若函数 \(f\) 在(a, b)上连续,则给定 \(\varepsilon  > 0\) ,对每一点 \(x \in  \left( {a,b}\right)\) ,都可确定正数 \({\delta }_{x}\) (它依赖于 \(\varepsilon\) 与 \(x\) ), 使得当 \({x}^{\prime } \in  U\left( {x;{\delta }_{x}}\right)\) 时,有 \(\left| {f\left( {x}^{\prime }\right)  - f\left( x\right) }\right|  < \varepsilon\) . 这样就得到一个开区间集
\[
H = \left\{  {\left( {x - {\delta }_{x},x + {\delta }_{x}}\right)  \mid  x \in  \left( {a,b}\right) }\right\}  ,
\]
它是区间(a, b)的一个无限开覆盖.

%定理 7.3
\begin{theorem}{}{the:}

\end{theorem}
 (海涅一博雷尔 (Heine-Borel) 有限覆盖定理) 设 \(H\) 为闭区间 \(\left\lbrack  {a,b}\right\rbrack\) 的一个 (无限) 开覆盖,则从 \(H\) 中可选出有限个开区间来覆盖 \(\left\lbrack  {a,b}\right\rbrack\) .

%证
\begin{proof}

\end{proof}
用反证法 假设定理的结论不成立,即不能用 \(H\) 中有限个开区间来覆盖 \(\left\lbrack  {a,b}\right\rbrack\) .

将 \(\left\lbrack  {a,b}\right\rbrack\) 等分为两个子区间,则其中至少有一个子区间不能用 \(H\) 中有限个开区间来覆盖. 记这个子区间为 \(\left\lbrack  {{a}_{1},{b}_{1}}\right\rbrack\) ,则 \(\left\lbrack  {{a}_{1},{b}_{1}}\right\rbrack   \subset  \left\lbrack  {a,b}\right\rbrack\) ,且 \({b}_{1} - {a}_{1} = \frac{1}{2}\left( {b - a}\right)\) .

再将 \(\left\lbrack  {{a}_{1},{b}_{1}}\right\rbrack\) 等分为两个子区间,同样,其中至少有一个子区间不能用 \(H\) 中有限个开区间来覆盖. 记这个子区间为 \(\left\lbrack  {{a}_{2},{b}_{2}}\right\rbrack\) ,则 \(\left\lbrack  {{a}_{2},{b}_{2}}\right\rbrack   \subset  \left\lbrack  {{a}_{1},{b}_{1}}\right\rbrack\) ,且 \({b}_{2} - {a}_{2} = \frac{1}{{2}^{2}}\left( {b - a}\right)\) .

重复上述步骤并不断地进行下去,则得到一个闭区间列 \(\left\{  \left\lbrack  {{a}_{n},{b}_{n}}\right\rbrack  \right\}\) ,它满足
\[
\left\lbrack  {{a}_{n},{b}_{n}}\right\rbrack   \supset  \left\lbrack  {{a}_{n + 1},{b}_{n + 1}}\right\rbrack  ,n = 1,2,\cdots ,
\]
\[
{b}_{n} - {a}_{n} = \frac{1}{{2}^{n}}\left( {b - a}\right)  \rightarrow  0\left( {n \rightarrow  \infty }\right) ,
\]
即 \(\left\{  \left\lbrack  {{a}_{n},{b}_{n}}\right\rbrack  \right\}\) 是区间套,且其中每一个闭区间都不能用 \(H\) 中有限个开区间来覆盖.

由区间套定理,存在唯一的一点 \(\xi  \in  \left\lbrack  {{a}_{n},{b}_{n}}\right\rbrack  ,n = 1,2,\cdots\) . 由于 \(H\) 是 \(\left\lbrack  {a,b}\right\rbrack\) 的一个开覆盖,故存在开区间 \(\left( {\alpha ,\beta }\right)  \in  H\) ,使 \(\xi  \in  \left( {\alpha ,\beta }\right)\) . 于是,由定理 7.1 推论,当 \(n\) 充分大时,有
\[
\left\lbrack  {{a}_{n},{b}_{n}}\right\rbrack   \subset  \left( {\alpha ,\beta }\right) .
\]
这表明 \(\left\lbrack  {{a}_{n},{b}_{n}}\right\rbrack\) 只需用 \(H\) 中的一个开区间 \(\left( {\alpha ,\beta }\right)\) 就能覆盖,与挑选 \(\left\lbrack  {{a}_{n},{b}_{n}}\right\rbrack\) 时的假设 “不能用 \(H\) 中有限个开区间来覆盖” 相矛盾. 从而证得必存在属于 \(H\) 的有限个开区间能覆盖 \(\left\lbrack  {a,b}\right\rbrack\) .

%注
\begin{remark}

\end{remark}
%定理 7.3
\begin{theorem}{}{the:}

\end{theorem}
 的结论只对闭区间 \(\left\lbrack  {a,b}\right\rbrack\) 成立,而对开区间则不一定成立. 例如,开区间集合 \(\left\{  \left( {\frac{1}{n + 1},1}\right) \right\}  \left( {n = 1,2,\cdots }\right)\) 构成了开区间(0,1)的一个开覆盖,但不能从中选出有限个开区间盖住(0,1).

%例 2
\begin{example}

\end{example}
用有限覆盖定理证明:闭区间上连续函数的有界性定理.

%证
\begin{proof}

\end{proof}
设 \(f\left( x\right)\) 在区间 \(\left\lbrack  {a,b}\right\rbrack\) 上连续. 根据连续函数的局部有界性定理,对于任意的 \({x}_{0} \in  \left\lbrack  {a,b}\right\rbrack\) ,存在正数 \({M}_{{x}_{0}}\) 以及正数 \({\delta }_{{x}_{0}}\) ,当 \(x \in  \left( {{x}_{0} - {\delta }_{{x}_{0}},{x}_{0} + {\delta }_{{x}_{0}}}\right)  \cap  \left\lbrack  {a,b}\right\rbrack\) 时有 \(\left| {f\left( x\right) }\right|  \leq\)  \({M}_{{x}_{0}}\) . 作开区间集
\[
H = \left\{  {\left( {x - {\delta }_{x},x + {\delta }_{x}}\right) \left| \right| f\left( x\right)  \mid   \leq  {M}_{x},x \in  \left\lbrack  {a,b}\right\rbrack  ,x \in  \left( {x - {\delta }_{x},x + {\delta }_{x}}\right)  \cap  \left\lbrack  {a,b}\right\rbrack  }\right\}  ,
\]
显然 \(H\) 覆盖了区间 \(\left\lbrack  {a,b}\right\rbrack\) . 根据有限覆盖定理,存在 \(H\) 中有限个开区间
\[
\left( {{x}_{1} - {\delta }_{{x}_{1}},{x}_{1} + {\delta }_{{x}_{1}}}\right) ,\left( {{x}_{2} - {\delta }_{{x}_{2}},{x}_{2} + {\delta }_{{x}_{2}}}\right) ,\cdots ,\left( {{x}_{n} - {\delta }_{{x}_{n}},{x}_{n} + {\delta }_{{x}_{n}}}\right) ,
\]
它们也覆盖了 \(\left\lbrack  {a,b}\right\rbrack\) . 令 \(M = \max \left\{  {{M}_{{x}_{1}},{M}_{{x}_{2}},\cdots ,{M}_{{x}_{n}}}\right\}\) ,那么对于任意的 \(x \in  \left\lbrack  {a,b}\right\rbrack\) ,存在 \(k\) , \(1 \leq  k \leq  n\) ,使得 \(x \in  \left( {{x}_{k} - {\delta }_{{x}_{k}},{x}_{k} + {\delta }_{{x}_{k}}}\right)\) ,并且有 \(\left| {f\left( x\right) }\right|  \leq  {M}_{{x}_{k}} \leq  M\) .

\subsection{实数完备性基本定理之间的等价性}

至此, 我们已经介绍了有关实数完备性的六个基本定理, 即

1. 确界原理 (定理 1.1);

2. 单调有界定理 (定理 2.9);

3. 区间套定理 (定理 7.1);

4. 有限覆盖定理 (定理 7.3);

5. 聚点定理 (定理 7.2) 和致密性定理 (定理 2.10);

6. 柯西收敛准则 (定理 2.11).

在本书中, 我们首先证明了确界原理, 由它证明了单调有界定理, 再用单调有界定理导出区间套定理, 最后用区间套定理分别证明余下的三个定理. 事实上, 在实数系中这六个命题是相互等价的, 即从其中任何一个命题都可推出其余的五个命题. 对此, 我们可按下列顺序给予证明:
\[
1 \Rightarrow  2 \Rightarrow  3 \Rightarrow  4 \Rightarrow  5 \Rightarrow  6 \Rightarrow  1\text{ . }
\]
其中 \(1 \Rightarrow  2,2 \Rightarrow  3\) 与 \(3 \Rightarrow  4\) 分别见定理 2.9,7.1 与 \({7.3};4 \Rightarrow  5\) 和 \(5 \Rightarrow  6\) 请读者作为练习自证 (见本节习题 8 和 9); 而 \(6 \Rightarrow  1\) 见下例.

%例 3
\begin{example}

\end{example}
用数列的柯西收敛准则证明确界原理.

%证
\begin{proof}

\end{proof}
设 \(S\) 为非空有上界数集. 由实数的阿基米德性,对任何正数 \(\alpha\) ,存在整数 \({k}_{\alpha }\) ,使得 \({\lambda }_{\alpha } = {k}_{\alpha }\alpha\) 为 \(S\) 的上界,而 \({\lambda }_{\alpha } - \alpha  = \left( {{k}_{\alpha } - 1}\right) \alpha\) 不是 \(S\) 的上界,即存在 \({\alpha }^{\prime } \in  S\) ,使得 \({\alpha }^{\prime } >\)  \(\left( {{k}_{\alpha } - 1}\right) \alpha\) .

分别取 \(\alpha  = \frac{1}{n},n = 1,2,\cdots\) ,则对每一个正整数 \(n\) ,存在相应的 \({\lambda }_{n}\) ,使得 \({\lambda }_{n}\) 为 \(S\) 的上界,而 \({\lambda }_{n} - \frac{1}{n}\) 不是 \(S\) 的上界,故存在 \({a}^{\prime } \in  S\) ,使得
\[
{a}^{\prime } > {\lambda }_{n} - \frac{1}{n}. \tag{6}
\]
又对正整数 \(m,{\lambda }_{m}\) 是 \(S\) 的上界,故有 \({\lambda }_{m} \geq  {a}^{\prime }\) . 结合 (6) 式得 \({\lambda }_{n} - {\lambda }_{m} < \frac{1}{n}\) ; 同理有 \({\lambda }_{m} - {\lambda }_{n} <\)  \(\frac{1}{m}\) . 从而得
\[
\left| {{\lambda }_{m} - {\lambda }_{n}}\right|  < \max \left\{  {\frac{1}{m},\frac{1}{n}}\right\}  .
\]
于是,对任给的 \(\varepsilon  > 0\) ,存在 \(N > 0\) ,使得当 \(m,n > N\) 时,有
\[
\left| {{\lambda }_{m} - {\lambda }_{n}}\right|  < \varepsilon
\]
由柯西收敛准则,数列 \(\left\{  {\lambda }_{n}\right\}\) 收敛. 记
\[
\mathop{\lim }\limits_{{n \rightarrow  \infty }}{\lambda }_{n} = \lambda  \tag{7}
\]
现在证明 \(\lambda\) 就是 \(S\) 的上确界. 首先,对任何 \(a \in  S\) 和正整数 \(n\) ,有 \(a \leq  {\lambda }_{n}\) ,由 (7) 式得 \(a \leq  \lambda\) ,即 \(\lambda\) 是 \(S\) 的一个上界. 其次,对任何 \(\delta  > 0\) ,由 \(\frac{1}{n} \rightarrow  0\left( {n \rightarrow  \infty }\right)\) 及 (7) 式,对充分大的 \(n\) ,同时有
\[
\frac{1}{n} < \frac{\delta }{2},{\lambda }_{n} > \lambda  - \frac{\delta }{2}.
\]
又因 \({\lambda }_{n} - \frac{1}{n}\) 不是 \(S\) 的上界,故存在 \({a}^{\prime } \in  S\) ,使得 \({a}^{\prime } > {\lambda }_{n} - \frac{1}{n}\) . 结合上式得
\[
{a}^{\prime } > \lambda  - \frac{\delta }{2} - \frac{\delta }{2} = \lambda  - \delta .
\]
这说明 \(\lambda\) 为 \(S\) 的上确界.

同理可证: 若 \(S\) 为非空有下界数集,则必存在下确界.

\subsection{习题 7.1}

1. 证明数集 \(\left\{  {{\left( -1\right) }^{n} + \frac{1}{n}}\right\}\) 有且只有两个聚点 \({\xi }_{1} =  - 1\) 和 \({\xi }_{2} = 1\) .

2. 证明: 任何有限数集都没有聚点.

3. 设 \(\left\{  \left( {{a}_{n},{b}_{n}}\right) \right\}\) 是一个严格开区间套,即满足
\[
{a}_{1} < {a}_{2} < \cdots  < {a}_{n} < {b}_{n} < \cdots  < {b}_{2} < {b}_{1},
\]
且 \(\mathop{\lim }\limits_{{n \rightarrow  \infty }}\left( {{b}_{n} - {a}_{n}}\right)  = 0\) . 证明: 存在唯一的一点 \(\xi\) ,使得
\[
{a}_{n} < \xi  < {b}_{n},n = 1,2,\cdots .
\]
4. 试举例说明:在有理数集上,确界原理、单调有界定理、聚点定理和柯西收敛准则一般都不能成立.

5. 设 \(H = \left\{  {\left( {\frac{1}{n + 2},\frac{1}{n}}\right)  \mid  n = 1,2,\cdots }\right\}\) . 问

(1)H 能否覆盖(0,1)？

(2)能否从 \(H\) 中选出有限个开区间覆盖(i) \(\left( {0,\frac{1}{2}}\right)\) ,(ii) \(\left( {\frac{1}{100},1}\right)\) ？

6. 证明: 闭区间 \(\left\lbrack  {a,b}\right\rbrack\) 的全体聚点的集合是 \(\left\lbrack  {a,b}\right\rbrack\) 本身.

7. 设 \(\left\{  {x}_{n}\right\}\) 为单调数列. 证明: 若 \(\left\{  {x}_{n}\right\}\) 存在聚点,则必是唯一的,且为 \(\left\{  {x}_{n}\right\}\) 的确界.

8. 试用有限覆盖定理证明聚点定理.

9. 试用聚点定理证明柯西收敛准则.

10. 用有限覆盖定理证明根的存在性定理.

11. 用有限覆盖定理证明连续函数的一致连续性定理.

\section{上极限和下极限}

%定义  1
\begin{definition}{}{def:}

\end{definition}
若数 \(a\) 的任一邻域含有数列 \(\left\{  {x}_{n}\right\}\) 中的无限多个项,则称 \(a\) 为数列 \(\left\{  {x}_{n}\right\}\) 的一个聚点\footnote{本节中同前面一样,不区分实数与数轴上的点,因此点列的聚点等同于数列的聚点. 数列或点列的聚点也称为极限点.}.

例如,数列 \(\left\{  {{\left( -1\right) }^{n}\frac{n}{n + 1}}\right\}\) 有聚点 -1 与 1 ; 数列 \(\left\{  {\sin \frac{n\pi }{4}}\right\}\) 有 \(- 1, - \frac{\sqrt{2}}{2},0,\frac{\sqrt{2}}{2}\) 和 1 五个聚点; 数列 \(\left\{  \frac{1}{n}\right\}\) 只有一个聚点 0 ; 常数列 \(\{ 1,1,\cdots ,1,\cdots \}\) 只有一个聚点 1 .

%注
\begin{remark}

\end{remark}
点列 (或数列) 的聚点定义与上一节中关于点集 (或数集) 的聚点定义是有区别的. 当把点列看作点集时, 点列中对应于相同数值的项, 只能作为一个点来看待. 如上述点列 \(\left\{  {\sin \frac{n\pi }{4}}\right\}\) 作为点集来看待时,它仅含有五个点,即
\[
\left\{  {\sin \frac{n\pi }{4}}\right\}   = \left\{  {-1, - \frac{\sqrt{2}}{2},0,\frac{\sqrt{2}}{2},1}\right\}  ,
\]
按点集聚点的定义, 这个有限集没有聚点. 然而, 我们在点列聚点的定义中只考虑项, 只要在一点的任意小邻域内聚集了无穷多个项 (不论其数值是否相同), 该点就称为点列的聚点. 所以, 点列的聚点实际上就是其收敛子列的极限.

%定理 7.4
\begin{theorem}{}{the:}

\end{theorem}
 有界点列 (数列) \(\left\{  {x}_{n}\right\}\) 至少有一个聚点,且存在最大聚点与最小聚点.

%证
\begin{proof}

\end{proof}
关于 \(\left\{  {x}_{n}\right\}\) 聚点存在性的证明,完全类似于定理 7.2 的证明方法,只需把那个证明中的“无限多个点”改为“无限多个项” 即可.

至于最大聚点的存在性,只需在定理 7.2 的证明过程中,当每次把区间 \(\left\lbrack  {{a}_{k - 1},{b}_{k - 1}}\right\rbrack\) 等分为两个子区间时,若右边一个含有 \(\left\{  {x}_{n}\right\}\) 中无穷多个项,则取它为 \(\left\lbrack  {{a}_{k},{b}_{k}}\right\rbrack\) ,否则取左边的子区间为 \(\left\lbrack  {{a}_{k},{b}_{k}}\right\rbrack\) . 这样的选取方法既保证了每次选出的 \(\left\lbrack  {{a}_{k},{b}_{k}}\right\rbrack\) 都含有 \(\left\{  {x}_{n}\right\}\) 中无限多个项,同时在 \(\left\lbrack  {{a}_{k},{b}_{k}}\right\rbrack\) 的右边却至多只有 \(\left\{  {x}_{n}\right\}\) 的有限个项,于是由区间套 \(\left\{  \left\lbrack  {{a}_{k},{b}_{k}}\right\rbrack  \right\}\) 所确定的点列 \(\left\{  {x}_{n}\right\}\) 的聚点 \(\xi\) 必是 \(\left\{  {x}_{n}\right\}\) 的最大聚点. 因若不然,设另有 \(\left\{  {x}_{n}\right\}\) 的聚点 \(\zeta  > \xi\) ,则令 \(\delta  = \frac{1}{3}\left( {\zeta  - \xi }\right)  > 0\) ,在 \(U\left( {\zeta ;\delta }\right)\) 内含有 \(\left\{  {x}_{n}\right\}\) 中无限多个项. 但当 \(n\) 充分大时, \(U\left( {\zeta ,\delta }\right)\) 将完全落在 \(\left\lbrack  {{a}_{n},{b}_{n}}\right\rbrack\) 的右边,这与区间列 \(\left\{  \left\lbrack  {{a}_{k},{b}_{k}}\right\rbrack  \right\}\) 的上述选取方法相矛盾. 所以 \(\xi\) 必为 \(\left\{  {x}_{n}\right\}\) 的最大聚点.

类似地, 只要把每次优先挑选右边一个子区间改为优先挑选左边一个, 就能证得最小聚点的存在性.

%定义  2
\begin{definition}{}{def:}

\end{definition}
有界数列 (点列) \(\left\{  {x}_{n}\right\}\) 的最大聚点 \(\bar{A}\) 与最小聚点 \(\underline{A}\) 分别称为 \(\left\{  {x}_{n}\right\}\) 的上极限与下极限, 记作
\[
\bar{A} = \mathop{\lim }\limits_{{n \rightarrow  \infty }}{x}_{n},\;\underline{A} = \mathop{\lim }\limits_{{n \rightarrow  \infty }}{x}_{n}.
\]


由定理 7.4 立刻可得:任何有界数列必存在上、下极限.

例 1
\[
\mathop{\lim }\limits_{{n \rightarrow  \infty }}{\left( -1\right) }^{n}\frac{n}{n + 1} = 1,\mathop{\lim }\limits_{{n \rightarrow  \infty }}{\left( -1\right) }^{n}\frac{n}{n + 1} =  - 1;
\]
\[
\mathop{\lim }\limits_{{n \rightarrow  \infty }}\sin \frac{n\pi }{4} = 1,\mathop{\lim }\limits_{{n \rightarrow  \infty }}\sin \frac{n\pi }{4} =  - 1;
\]
\[
\mathop{\lim }\limits_{{n \rightarrow  \infty }}\frac{1}{n} = \mathop{\lim }\limits_{{n \rightarrow  \infty }}\frac{1}{n} = 0.
\]
%定理 7.5
\begin{theorem}{}{the:}

\end{theorem}
 对任何有界数列 \(\left\{  {x}_{n}\right\}\) ,有
\[
\mathop{\lim }\limits_{{n \rightarrow  \infty }}{x}_{n} \leq  \mathop{\lim }\limits_{{n \rightarrow  \infty }}{x}_{n}.
\]
定理 \({7.6}\mathop{\lim }\limits_{{n \rightarrow  \infty }}{x}_{n} = A\) 的充要条件是 \(\mathop{\lim }\limits_{{n \rightarrow  \infty }}{x}_{n} = \mathop{\lim }\limits_{{n \rightarrow  \infty }}{x}_{n} = A\) .

以上两个定理的证明由定理 7.4 与定义 2 立即可得.

%定理 7.7
\begin{theorem}{}{the:}

\end{theorem}
 设 \(\left\{  {x}_{n}\right\}\) 为有界数列.

(1) \(\bar{A}\) 为 \(\left\{  {x}_{n}\right\}\) 上极限的充要条件是:任给 \(\varepsilon  > 0\) ,

(i) 存在 \(N > 0\) ,使得当 \(n > N\) 时,有 \({x}_{n} < \bar{A} + \varepsilon\) ;

(ii) 存在子列 \(\left\{  {x}_{{n}_{k}}\right\}  ,{x}_{{n}_{k}} > \bar{A} - \varepsilon ,k = 1,2,\cdots\) .

(2) \(\underline{A}\) 为 \(\left\{  {x}_{n}\right\}\) 下极限的充要条件是:任给 \(\varepsilon  > 0\) ,

(i) 存在 \(N > 0\) ,使得当 \(n > N\) 时,有 \({x}_{n} > A - \varepsilon\) ;

(ii) 存在子列 \(\left\{  {x}_{{n}_{k}}\right\}  ,{x}_{{n}_{k}} < \underline{A} + \varepsilon ,k = 1,2,\cdots\) .

%证
\begin{proof}

\end{proof}
(1) 必要性 因 \(\bar{A}\) 是 \(\left\{  {x}_{n}\right\}\) 的聚点,故对任给的 \(\varepsilon  > 0,U\left( {\bar{A};\varepsilon }\right)\) 含有 \(\left\{  {x}_{n}\right\}\) 中无穷多项,设为 \(\left\{  {x}_{{n}_{k}}\right\}\) ,则有 \({x}_{{n}_{k}} > \bar{A} - \varepsilon ,k = 1,2,\cdots\) .

又因 \(\bar{A}\) 是 \(\left\{  {x}_{n}\right\}\) 的最大聚点,故在 \(\bar{A} + \varepsilon\) 的右边至多只有 \(\left\{  {x}_{n}\right\}\) 的有限个项,设此有限项的最大下标为 \(N\) ,则当 \(n > N\) 时,有 \({x}_{n} < \bar{A} + \varepsilon\) .

充分性 任给 \(\varepsilon  > 0\) ,由条件 (i) 和 (ii) 易见, \(U\left( {\bar{A};\varepsilon }\right)\) 含有 \(\left\{  {x}_{n}\right\}\) 中无穷多个项,故 \(\bar{A}\) 是 \(\left\{  {x}_{n}\right\}\) 的一个聚点.

又设 \(\alpha  > \bar{A}\) . 记 \(\varepsilon  = \frac{1}{2}\left( {\alpha  - \bar{A}}\right)\) ,则由条件 (i) 易见 \(U\left( {\alpha ;\varepsilon }\right)\) 内至多只有 \(\left\{  {x}_{n}\right\}\) 中有限个项,故 \(\alpha\) 不是 \(\left\{  {x}_{n}\right\}\) 的聚点. 所以 \(\bar{A}\) 是 \(\left\{  {x}_{n}\right\}\) 的最大聚点.

(2)类似地证明.

%定理 7.7
\begin{theorem}{}{the:}

\end{theorem}
 的另一种形式如下:

定理 \({7.7}^{\prime }\) 设 \(\left\{  {x}_{n}\right\}\) 为有界数列.

(1) \(\bar{A}\) 为 \(\left\{  {x}_{n}\right\}\) 上极限的充要条件是: 对任何 \(\alpha  > \bar{A},\left\{  {x}_{n}\right\}\) 中大于 \(\alpha\) 的项至多有限个; 对任何 \(\beta  < \bar{A},\left\{  {x}_{n}\right\}\) 中大于 \(\beta\) 的项有无限多个.

(2) \(\underline{A}\) 为 \(\left\{  {x}_{n}\right\}\) 下极限的充要条件是:对任何 \(\beta  < \underline{A},\left\{  {x}_{n}\right\}\) 中小于 \(\beta\) 的项至多有限个； 对任何 \(\alpha  > \underline{A},\left\{  {x}_{n}\right\}\) 中小于 \(\alpha\) 的项有无限多个.

%定理 7.8
\begin{theorem}{}{the:}

\end{theorem}
(上、下极限的保不等式性) 设有界数列 \(\left\{  {a}_{n}\right\}  ,\left\{  {b}_{n}\right\}\) 满足: 存在 \({N}_{0} > 0\) ,当 \(n > {N}_{0}\) 时,有 \({a}_{n} \leq  {b}_{n}\) ,则
\[
\mathop{\lim }\limits_{{n \rightarrow  \infty }}{a}_{n} \leq  \mathop{\lim }\limits_{{n \rightarrow  \infty }}{b}_{n},\mathop{\lim }\limits_{{n \rightarrow  \infty }}{a}_{n} \leq  \mathop{\lim }\limits_{{n \rightarrow  \infty }}{b}_{n}.
\]
特别地,若 \(\alpha ,\beta\) 为常数,又存在 \({N}_{0} > 0\) ,当 \(n > {N}_{0}\) 时,有 \(\alpha  \leq  {a}_{n} \leq  \beta\) ,则
\[
\alpha  \leq  \mathop{\lim }\limits_{\overline{n \rightarrow  \infty }}{a}_{n} \leq  \mathop{\lim }\limits_{{n \rightarrow  \infty }}{a}_{n} \leq  \beta .
\]
这个定理的证明留给读者.

%例 2
\begin{example}

\end{example}
设 \(\left\{  {a}_{n}\right\}  ,\left\{  {b}_{n}\right\}\) 为有界数列. 证明
\[
\mathop{\lim }\limits_{{n \rightarrow  \infty }}\left( {{a}_{n} + {b}_{n}}\right)  \leq  \mathop{\lim }\limits_{{n \rightarrow  \infty }}{a}_{n} + \mathop{\lim }\limits_{{n \rightarrow  \infty }}{b}_{n}. \tag{1}
\]
特别地,若 \(\mathop{\lim }\limits_{{n \rightarrow  \infty }}{a}_{n}\) 存在,则
\[
\mathop{\lim }\limits_{{n \rightarrow  \infty }}\left( {{a}_{n} + {b}_{n}}\right)  = \mathop{\lim }\limits_{{n \rightarrow  \infty }}{a}_{n} + \mathop{\lim }\limits_{{n \rightarrow  \infty }}{b}_{n} = \mathop{\lim }\limits_{{n \rightarrow  \infty }}{a}_{n} + \mathop{\lim }\limits_{{n \rightarrow  \infty }}{b}_{n}. \tag{2}
\]
%证
\begin{proof}

\end{proof}
设 \(\mathop{\lim }\limits_{{n \rightarrow  \infty }}{a}_{n} = A,\mathop{\lim }\limits_{{n \rightarrow  \infty }}{b}_{n} = B\) . 由定理 7.7,对任给的 \(\varepsilon  > 0\) ,存在 \(N > 0\) ,当 \(n > N\) 时,有
\[
{a}_{n} < A + \frac{\varepsilon }{2},{b}_{n} < B + \frac{\varepsilon }{2} \Rightarrow  {a}_{n} + {b}_{n} < A + B + \varepsilon .
\]
再利用上极限的保不等式性 (定理 7.8) 得
\[
\mathop{\lim }\limits_{{n \rightarrow  \infty }}\left( {{a}_{n} + {b}_{n}}\right)  \leq  A + B + \varepsilon .
\]
故由 \(\varepsilon\) 的任意性得 \(\mathop{\lim }\limits_{{n \rightarrow  \infty }}\left( {{a}_{n} + {b}_{n}}\right)  \leq  A + B\) ,即 (1) 式成立.

若 \(\mathop{\lim }\limits_{{n \rightarrow  \infty }}{a}_{n} = \mathop{\lim }\limits_{{n \rightarrow  \infty }}{a}_{n} = A\) (即极限存在),由 (1) 式得
\[
\mathop{\lim }\limits_{{n \rightarrow  \infty }}{b}_{n} = \mathop{\lim }\limits_{{n \rightarrow  \infty }}\left\lbrack  {\left( {{a}_{n} + {b}_{n}}\right)  - {a}_{n}}\right\rbrack   \leq  \mathop{\lim }\limits_{{n \rightarrow  \infty }}\left( {{a}_{n} + {b}_{n}}\right)  + \overline{\lim }\left( {-{a}_{n}}\right)
\]
\[
= \mathop{\lim }\limits_{{n \rightarrow  \infty }}\left( {{a}_{n} + {b}_{n}}\right)  - A\text{ . }
\]
故
\[
\mathop{\lim }\limits_{{n \rightarrow  \infty }}\left( {{a}_{n} + {b}_{n}}\right)  \geq  \mathop{\lim }\limits_{{n \rightarrow  \infty }}{a}_{n} + \mathop{\lim }\limits_{{n \rightarrow  \infty }}{b}_{n},
\]
结合 (1) 式可知 (2) 式成立.

%注
\begin{remark}

\end{remark}
(1) 式有可能成立严格的不等式. 例如,设 \({a}_{n} = {\left( -1\right) }^{n},{b}_{n} = {\left( -1\right) }^{n + 1}\) ,则易见 (1) 式左边等于 0 , 右边等于 2 .

%定理 7.9
\begin{theorem}{}{the:}

\end{theorem}
 设 \(\left\{  {x}_{n}\right\}\) 为有界数列.

(1) \(\bar{A}\) 为 \(\left\{  {x}_{n}\right\}\) 上极限的充要条件是
\[
\bar{A} = \mathop{\lim }\limits_{{n \rightarrow  \infty }}\mathop{\sup }\limits_{{k \geq  n}}\left\{  {x}_{k}\right\}  ; \tag{3}
\]
(2) \(\underline{A}\) 为 \(\left\{  {x}_{n}\right\}\) 下极限的充要条件是
\[
\underline{A} = \mathop{\lim }\limits_{{n \rightarrow  \infty }}\mathop{\inf }\limits_{{k \geq  n}}\left\{  {x}_{k}\right\}  . \tag{4}
\]
做过第二章 § 3 习题 12 的读者, 对这个定理应该不会感到陌生, 并能自行写出其证明. 有些教科书上也用 (3)、(4) 分别作为有界数列 \(\left\{  {x}_{n}\right\}\) 上、下极限的定义.

若定义 1 中的 \(a\) 可允许是非正常点 + ∞ 或 - ∞,则定理 7.4 可相应地扩充为: 任一点列 \(\left\{  {x}_{n}\right\}\) 至少有一个聚点,且存在最大聚点与最小聚点.

不难证明:无上 (下) 界点列的最大 (小) 聚点为 \(+ \infty \left( {-\infty }\right)\) . 于是,无上 (下) 界点列有非正常上 (下) 极限 \(+ \infty \left( {-\infty }\right)\) . 例如,
\[
\mathop{\lim }\limits_{{n \rightarrow  \infty }}\left\lbrack  {{\left( -1\right) }^{n} + 1}\right\rbrack  n =  + \infty ,\;\mathop{\lim }\limits_{{n \rightarrow  \infty }}\left\lbrack  {{\left( -1\right) }^{n} + 1}\right\rbrack  n = 0;
\]
\[
\mathop{\lim }\limits_{{n \rightarrow  \infty }}{\left( -1\right) }^{n}n =  + \infty ,\mathop{\lim }\limits_{{n \rightarrow  \infty }}{\left( -1\right) }^{n}n =  - \infty .
\]
%注
\begin{remark}

\end{remark}
对于非正常上、下极限,上述定理 7.5 至 7.9 也成立(其中定理 7.7 应作相应地修改. 例如, \(\mathop{\lim }\limits_{{n \rightarrow  \infty }}{x}_{n} =  + \infty\) 的充要条件是
\[
\mathop{\lim }\limits_{{n \rightarrow  \infty }}{x}_{n} = \mathop{\lim }\limits_{{n \rightarrow  \infty }}{x}_{n} =  + \infty )\text{ . }
\]
\subsection{习题 7.2}

1. 求以下数列的上、下极限:

(1) \(\left\{  {1 + {\left( -1\right) }^{n}}\right\}\) ;

(2) \(\left\{  {{\left( -1\right) }^{n}\frac{n}{{2n} + 1}}\right\}\) ;

(3) \(\{ {{2n} + 1}\}\) ;

(4) \(\left\{  {\frac{2n}{n + 1}\sin \frac{n\pi }{4}}\right\}\) ;

(5) \(\left\{  {\frac{{n}^{2} + 1}{n}\sin \frac{\pi }{n}}\right\}\) ;

(6) \(\left\{  \sqrt[n]{\left| \cos \frac{n\pi }{3}\right| }\right\}\) .

2. 设 \(\left\{  {a}_{n}\right\}  ,\left\{  {b}_{n}\right\}\) 为有界数列,证明:

(1) \(\mathop{\lim }\limits_{\overline{n \rightarrow  \infty }}{a}_{n} =  - \mathop{\lim }\limits_{{n \rightarrow  \infty }}\left( {-{a}_{n}}\right)\) ;

(2) \(\mathop{\lim }\limits_{{n \rightarrow  \infty }}{a}_{n} + \mathop{\lim }\limits_{{n \rightarrow  \infty }}{b}_{n} \leq  \mathop{\lim }\limits_{{n \rightarrow  \infty }}\left( {{a}_{n} + {b}_{n}}\right)\) ;

(3)若 \({a}_{n} > 0,{b}_{n} > 0\left( {n = 1,2,\cdots }\right)\) ,则
\[
\mathop{\lim }\limits_{{n \rightarrow  \infty }}{a}_{n}\mathop{\lim }\limits_{{n \rightarrow  \infty }}{b}_{n} \leq  \mathop{\lim }\limits_{{n \rightarrow  \infty }}{a}_{n}{b}_{n},\mathop{\lim }\limits_{{n \rightarrow  \infty }}{a}_{n}\mathop{\lim }\limits_{{n \rightarrow  \infty }}{b}_{n} \geq  \mathop{\lim }\limits_{{n \rightarrow  \infty }}{a}_{n}{b}_{n};
\]
(4)若 \({a}_{n} > 0,\mathop{\lim }\limits_{{n \rightarrow  \infty }}{a}_{n} > 0\) ,则
\[
\mathop{\lim }\limits_{{n \rightarrow  \infty }}\frac{1}{{a}_{n}} = \frac{1}{\mathop{\lim }\limits_{{n \rightarrow  \infty }}{a}_{n}}.
\]
3. 证明: 若 \(\left\{  {a}_{n}\right\}\) 为递增数列,则 \(\mathop{\lim }\limits_{{n \rightarrow  \infty }}{a}_{n} = \mathop{\lim }\limits_{{n \rightarrow  \infty }}{a}_{n}\) .

4. 证明: 若 \({a}_{n} > 0\left( {n = 1,2,\cdots }\right)\) 且 \(\mathop{\lim }\limits_{{n \rightarrow  \infty }}{a}_{n} \cdot  \mathop{\lim }\limits_{{n \rightarrow  \infty }}\frac{1}{{a}_{n}} = 1\) ,则数列 \(\left\{  {a}_{n}\right\}\) 收敛.

5. 证明定理 7.8 .

6. 证明定理 7.9.

\section{第七章总练习题}

1. 设 \({E}^{\prime }\) 是集合 \(E\) 的全体聚点所成的点集, \({x}_{0}\) 是 \({E}^{\prime }\) 的一个聚点. 试证: \({x}_{0} \in  {E}^{\prime }\) .

2. 用确界原理证明有限覆盖定理.

3. 设 \(\mathop{\lim }\limits_{{n \rightarrow  \infty }}{x}_{n} = A < B = \mathop{\lim }\limits_{{n \rightarrow  \infty }}{x}_{n},\mathop{\lim }\limits_{{n \rightarrow  \infty }}\left( {{x}_{n + 1} - {x}_{n}}\right)  = 0\) . 试证: 数列 \(\left\{  {x}_{n}\right\}\) 的聚点全体恰为闭区间 \(\left\lbrack  {A,B}\right\rbrack\) .

